% =====================================================================
%  chapters/minggu-01.tex
%  Minggu 1 — Pengantar Web \& Setup Lingkungan Kerja
% =====================================================================

\chapter{Pengantar Web \& Setup Lingkungan Kerja}
\label{chap:minggu-01}

% ---------------------------------------------------------------------
% 1. CAPAIAN PEMBELAJARAN
% ---------------------------------------------------------------------
\begin{capaian}
Setelah menyelesaikan bab ini, mahasiswa diharapkan mampu:
\begin{itemize}
  \item menjelaskan alur kerja \emph{design-to-code} dalam pengembangan web;
  \item menginstal dan mengonfigurasi VS Code beserta ekstensi yang diperlukan;
  \item membuat akun GitHub dan melakukan \emph{push} pertama ke \emph{repository};
  \item membuka file HTML di peramban menggunakan \emph{Live Server}.
\end{itemize}
\end{capaian}

% ---------------------------------------------------------------------
% 2. TUJUAN PRAKTIK
% ---------------------------------------------------------------------
\begin{tujuanpraktik}
Pada sesi praktik ini, mahasiswa akan:
\begin{itemize}
  \item menginstal VS Code dan ekstensi pendukung;
  \item membuat file HTML pertama dan melihat hasilnya di peramban;
  \item membuat akun GitHub serta \emph{repository} pertama;
  \item melakukan \emph{push} kode ke GitHub.
\end{itemize}
\end{tujuanpraktik}

% ---------------------------------------------------------------------
% 3. PENGANTAR: APA ITU WEB?
% ---------------------------------------------------------------------
\section{Pengenalan Web}
\label{sec:m01-pengenalan}

Website adalah kumpulan halaman yang diakses melalui peramban (\emph{browser}) seperti Chrome, Firefox, atau Edge. Setiap halaman web terdiri dari tiga lapis utama:

\begin{table}[ht]
\centering
\caption{Tiga lapis utama halaman web}
\label{tab:m01-lapis}
\begin{tabular}{llp{7cm}}
\toprule
\textbf{Lapis} & \textbf{Bahasa} & \textbf{Fungsi} \\
\midrule
Struktur & HTML & ``Tulang'' — menentukan elemen apa yang ada di halaman \\
Tampilan & CSS & ``Kulit'' — menentukan warna, ukuran, posisi, font \\
Interaksi & JavaScript & ``Otot'' — menentukan apa yang terjadi saat diklik/di-scroll \\
\bottomrule
\end{tabular}
\end{table}

\begin{catatan}
Mahasiswa tidak perlu menghafal ketiga lapis ini sekarang. Cukup ketahui bahwa ada tiga lapis; di minggu-minggu berikutnya kita akan pelajari satu per satu.
\end{catatan}

% ---------------------------------------------------------------------
% 4. ALUR KERJA DESIGN-TO-CODE
% ---------------------------------------------------------------------
\section{Alur Kerja \emph{Design-to-Code}}
\label{sec:m01-alur}

Alur kerja pengembangan web mengikuti urutan berikut:

\begin{enumerate}
  \item \textbf{Ide / Kebutuhan} — menentukan tujuan dan isi halaman.
  \item \textbf{Figma} — mendesain tampilan visual.
  \item \textbf{VS Code} — menulis kode HTML, CSS, dan JavaScript.
  \item \textbf{Peramban} — melihat dan menguji hasil.
  \item \textbf{GitHub} — menyimpan kode agar bisa diakses kapan saja.
\end{enumerate}

Intinya: desain di Figma\,$\rightarrow$\,implementasi di VS Code\,$\rightarrow$\,lihat hasil di peramban\,$\rightarrow$\,simpan di GitHub.

% ---------------------------------------------------------------------
% 5. LIVE CODING: INSTALASI VS CODE \& EKSTENSI
% ---------------------------------------------------------------------
\section{Instalasi VS Code \& Ekstensi}
\label{sec:m01-instalasi}

\subsection{Download \& Instal VS Code}

\begin{langkah}
  \item Buka peramban, kunjungi \texttt{https://code.visualstudio.com/}.
  \item Klik tombol \textbf{Download for Windows} (atau sesuai OS).
  \item Jalankan installer, klik \textbf{Next} di setiap langkah, pastikan opsi berikut diceklis:
  \begin{itemize}
    \item Add ``Open with Code'' action to Windows Explorer file context menu
    \item Add ``Open with Code'' action to Windows Explorer directory context menu
    \item Register Code as an editor for supported file types
  \end{itemize}
  \item Klik \textbf{Install}, tunggu selesai, lalu klik \textbf{Finish}.
\end{langkah}

\subsection{Instal Ekstensi Penting}

Buka VS Code, tekan \textbf{Ctrl + Shift + X} untuk membuka panel \emph{Extensions}. Cari dan instal ekstensi berikut:

\begin{table}[ht]
\centering
\caption{Ekstensi VS Code yang diperlukan}
\label{tab:m01-ekstensi}
\begin{tabular}{llp{6cm}}
\toprule
\textbf{No.} & \textbf{Nama Ekstensi} & \textbf{Fungsi} \\
\midrule
1 & Live Server (Ritwick Dey) & Preview halaman HTML otomatis di peramban saat file disimpan \\
2 & HTML CSS Support & Autocomplete kelas CSS di file HTML \\
3 & Prettier -- Code formatter & Format kode otomatis agar rapi \\
4 & VS Code Intellicode & Saran kode cerdas dari AI \\
\bottomrule
\end{tabular}
\end{table}

\subsection{Verifikasi Instalasi}

\begin{langkah}
  \item Buka VS Code.
  \item Buat folder baru: \textbf{File\,$\rightarrow$\,Open Folder\,$\rightarrow$\,New Folder}, beri nama \texttt{latihan-web-1}.
  \item Di dalam folder itu, buat file baru bernama \texttt{index.html}.
  \item Ketik \texttt{!} lalu tekan \textbf{Tab} — VS Code akan mengisi template HTML secara otomatis (\emph{snippets} Emmet).
  \item Ubah isi \texttt{<title>} menjadi \texttt{Halo Dunia}.
  \item Tambahkan di dalam \texttt{<body>}:
\end{langkah}

\begin{lstlisting}[language=HTML, caption={Kode ``Halo Dunia''}, label={lst:m01-halo}]
<h1>Halo, Dunia!</h1>
<p>Ini halaman web pertamaku.</p>
\end{lstlisting}

\begin{langkah}
  \item Klik kanan di editor $\rightarrow$ \textbf{Open with Live Server}.
  \item Peramban akan terbuka secara otomatis menampilkan teks ``Halo, Dunia!''.
\end{langkah}

\begin{figure}[htbp]
  \centering
  \includegraphics[width=0.85\linewidth]{minggu-01/vs-code-index.html.png}
  \caption{VS Code dengan file \texttt{index.html} terbuka.}
  \label{fig:vs-code-index}
\end{figure}

\begin{figure}[htbp]
  \centering
  \includegraphics[width=0.85\linewidth]{minggu-01/browser-halo-dunia.png}
  \caption{Peramban menampilkan hasil ``Halo, Dunia!''}
  \label{fig:browser-halo-dunia}
\end{figure}

% ---------------------------------------------------------------------
% 6. LIVE CODING: AKUN GITHUB \& REPOSITORY PERTAMA
% ---------------------------------------------------------------------
\section{Akun GitHub \& Repository Pertama}
\label{sec:m01-github}

\subsection{Buat Akun GitHub}

\begin{langkah}
  \item Buka \texttt{https://github.com}.
  \item Klik \textbf{Sign up}.
  \item Isi email, password, dan username (gunakan format: \texttt{nama-nim}, misal: \texttt{budi-2301001}).
  \item Selesaikan verifikasi.
\end{langkah}

\begin{figure}[htbp]
  \centering
  \includegraphics[width=0.85\linewidth]{minggu-01/dashboard-github.png}
  \caption{Dashboard GitHub setelah berhasil login.}
  \label{fig:dashboard-github}
\end{figure}

\subsection{Buat Repository Baru}

\begin{langkah}
  \item Di dashboard GitHub, klik tombol \textbf{+} (pojok kanan atas) $\rightarrow$ \textbf{New repository}.
  \item Isi form berikut:
  \begin{itemize}
    \item \textbf{Repository name:} \texttt{latihan-web-1}
    \item \textbf{Description:} \texttt{Latihan pertama --- Halo Dunia}
    \item Pilih \textbf{Public}
    \item Centang \textbf{Add a README file}
  \end{itemize}
  \item Klik \textbf{Create repository}.
\end{langkah}

\subsection{Hubungkan Folder Lokal ke GitHub (\emph{Push} Pertama)}

\begin{peringatan}
Pastikan Git sudah terinstal di laptop. Jika belum, unduh dari \texttt{https://git-scm.com/download/win} dan instal dengan pengaturan default.
\end{peringatan}

Buka \textbf{Terminal} di VS Code (\textbf{Ctrl + \$} atau menu \textbf{Terminal\,$\rightarrow$\,New Terminal}), lalu ketik perintah berikut satu per satu:

\begin{lstlisting}[language=bash, caption={Perintah push pertama ke GitHub}, label={lst:m01-git-push}]
# 1. Inisialisasi Git di folder lokal
git init

# 2. Hubungkan ke remote repository
#    (ganti NAMA_USER dengan username GitHub-mu)
git remote add origin https://github.com/NAMA_USER/latihan-web-1.git

# 3. Tambahkan semua file untuk di-commit
git add .

# 4. Buat commit pertama
git commit -m "Commit pertama: Halo Dunia"

# 5. Push ke GitHub
git branch -M main
git push -u origin main
\end{lstlisting}

\begin{catatan}
Jika Git meminta \emph{credential} (username \& password), gunakan \textbf{Personal Access Token} sebagai password:
\begin{enumerate}
  \item Buka \texttt{https://github.com/settings/tokens}.
  \item Klik \textbf{Generate new token (classic)}.
  \item Beri nama token, centang \textbf{repo}, lalu klik \textbf{Generate}.
  \item Salin token dan gunakan sebagai password saat Git meminta.
\end{enumerate}
\end{catatan}

\subsection{Verifikasi}

Buka \texttt{https://github.com/NAMA\_USER/latihan-web-1} di peramban. Pastikan file \texttt{index.html} terlihat di daftar file.

\begin{figure}[htbp]
  \centering
  \includegraphics[width=0.85\linewidth]{minggu-01/repository-github.png}
  \caption{Halaman repository GitHub \texttt{latihan-web-1}.}
  \label{fig:repository-github}
\end{figure}

% ---------------------------------------------------------------------
% 7. LATIHAN MANDIRI
% ---------------------------------------------------------------------
\section{Latihan Mandiri}
\label{sec:m01-latihan}

\begin{latihan}[Perbarui Halaman]
Buka file \texttt{index.html} di VS Code. Tambahkan elemen berikut di bawah \texttt{<p>} yang sudah ada:
\begin{lstlisting}[language=HTML]
<h2>Tentang Aku</h2>
<p>Nama: [nama lengkap kamu]</p>
<p>NIM: [NIM kamu]</p>
<p>Prodi: D3 Informatika</p>
\end{lstlisting}
Simpan file (\textbf{Ctrl + S}), periksa peramban, lalu \emph{commit} dan \emph{push} perubahan ke GitHub.
\end{latihan}

\begin{latihan}[Eksperimen \emph{Live Server}]
\begin{enumerate}
  \item Matikan \emph{Live Server} (klik \textbf{Port: 5500\,$\rightarrow$\,Stop} di pojok kanan bawah VS Code).
  \item Buka file \texttt{index.html} langsung dari File Explorer (klik ganda).
  \item Perhatikan URL di peramban --- apakah berbeda dari URL \emph{Live Server}? (\texttt{file:///...} vs \texttt{127.0.0.1:5500})
  \item Ubah salah satu teks di VS Code $\rightarrow$ simpan $\rightarrow$ muat ulang peramban secara manual.
  \item Aktifkan kembali \emph{Live Server}. Ubah teks lagi $\rightarrow$ simpan $\rightarrow$ perhatikan apakah peramban memperbarui diri secara otomatis.
\end{enumerate}
\end{latihan}

\begin{figure}[htbp]
  \centering
  \includegraphics[width=0.85\linewidth]{minggu-01/perbandingan-url.png}
  {Perbandingan URL: \emph{Live Server} vs file langsung.}
  {Kiri: \texttt{127.0.0.1:5500/index.html}; kanan: \texttt{file:///C:/.../index.html}.}
\end{figure}

% ---------------------------------------------------------------------
% 8. TROUBLESHOOTING
% ---------------------------------------------------------------------
\section{Troubleshooting}
\label{sec:m01-troubleshoot}

\begin{table}[ht]
\centering
\caption{Masalah umum dan solusinya}
\label{tab:m01-troubleshoot}
\small
\begin{tabular}{p{4.5cm}p{8cm}}
\toprule
\textbf{Masalah} & \textbf{Solusi} \\
\midrule
\emph{Live Server} tidak muncul di menu kanan-klik & Pastikan ekstensi ``Live Server'' sudah terinstal dan VS Code di-\emph{restart} \\
Peramban menampilkan \texttt{This site can't be reached} & Pastikan \emph{Live Server} sudah berjalan (lihat \emph{status bar} pojok kanan bawah VS Code) \\
Git menolak \emph{push} karena ``repository not found'' & Periksa URL \emph{remote}: \texttt{git remote -v}, lalu perbaiki dengan \texttt{git remote set-url origin URL\_BARU} \\
Git meminta \emph{password} tapi tidak bisa login & Gunakan \textbf{Personal Access Token} (bukan \emph{password} GitHub biasa) \\
VS Code menampilkan error ``command not found: git'' & Instal Git dari \texttt{https://git-scm.com/download/win}, lalu \emph{restart} VS Code \\
File \texttt{index.html} tidak muncul di GitHub & Pastikan sudah menjalankan \texttt{git add .}, \texttt{git commit}, dan \texttt{git push} \\
\bottomrule
\end{tabular}
\end{table}

% ---------------------------------------------------------------------
% 9. RANGKUMAN
% ---------------------------------------------------------------------
\section{Rangkuman}
\label{sec:m01-rangkuman}

\begin{itemize}
  \item Alur kerja \emph{design-to-code}: Figma $\rightarrow$ VS Code $\rightarrow$ Peramban $\rightarrow$ GitHub.
  \item VS Code + ekstensi (\emph{Live Server}, Prettier, dll.) sudah terinstal.
  \item File HTML pertama berhasil dibuat dan dilihat di peramban.
  \item Akun GitHub dan \emph{repository} pertama sudah dibuat.
  \item \emph{Push} pertama ke GitHub berhasil dilakukan.
  \item \emph{Live Server} memberikan \emph{preview} otomatis tanpa \emph{reload}.
\end{itemize}
